A system that acts on the world and is then scored on what it observes can write
part of its own objective. In patrol allocation, patrolling a district raises
the fraction of its crime that is recorded, and recorded crime is both the input
and the success measure of the next allocation. We study what happens when such
a \emph{self-contaminated proxy} is optimised, using an agent-based crime
simulator that is never seeded with historical crime, so latent crime is known
everywhere. Policies set the parameters of an allocation rule and are rewarded
for covering \emph{recorded} crime; each is scored against a need-proportional
reference on the same simulator seeds. A reinforcement-learning agent (PPO)
places \nPPOx$\times$ the reference's patrol disparity on the most disadvantaged
districts (95\% CI \nPPOxLo--\nPPOxHi) while covering less true crime than the
reference, and a non-learning search that maximises the same proxy is more
skewed still (\nBRRx$\times$): the bias belongs to the objective, and better
optimisation makes it worse. A mean-field model of the loop explains why.
Patrol-dependent recording makes the recorded signal rise with patrol, which
amplifies the tilt of any allocation rule and keeps the proxy reward increasing
with concentration after true-crime coverage has peaked; the model reproduces
the simulator's disparities to within \nThDispErr\% on average. Removing the
dependence of detection on patrol cuts PPO's excess from \nEfiveBase$\times$ to
\nEfiveConst$\times$. Parity penalties overcorrect and under-patrol the
districts with the most crime, whereas a penalty targeted at the
need-proportional reference lands closest to it while raising true-crime
coverage.
